\begin{table}[H]
\centering
\scriptsize
\renewcommand{\arraystretch}{1.15}
\setlength{\tabcolsep}{4pt}
\begin{tabular}{|l|r|r|r|r|r|r|r|r|r|}
\hline
\multicolumn{10}{|c|}{\textbf{DEPRECIACIÓN DE ACTIVOS TANGIBLES}} \\ \hline
\textbf{Maquinaria y Equipo} & \textbf{Cantidad} & \textbf{C/unitario} & \textbf{costo total} & \textbf{1\%} & \textbf{Valor depreciar} & \textbf{Vida útil} & \textbf{Dprc. Anual} & \textbf{Dprc. Acum} & \textbf{Vr. Salvamiento} \\ \hline
Equipo y Mobiliario Administración & 3 & 1,500 & 4,500 & 45 & 4,455 & 6 & 742.5 & 7,425 & -2,925 \\ \hline
Equipo y Mobiliario Ventas & 2 & 1,500 & 3,000 & 30 & 2,970 & 6 & 495 & 4,950 & -1,950 \\ \hline
Equipo y Mobiliario Desarrollo & 5 & 3,000 & 15,000 & 150 & 14,850 & 6 & 2,475 & 24,750 & -9,750 \\ \hline
\textbf{total} &  &  & \textbf{22,500.00} & \textbf{225.00} & \textbf{22,275.00} &  & \textbf{3,712.50} & \textbf{37,125.00} & \textbf{-14,625.00} \\ \hline
\end{tabular}

\vspace{0.2cm}
\textbf{Figura 10.18.} Depreciación de Activos tangibles \textbf{[Elaboración Propia]}
\end{table}

\subsection{Plan de Inversión y Capital de Trabajo}

Para realizar el plan de inversión tomamos en cuenta tres aspectos fundamentales: el total de nuestros activos fijos (mobiliario y equipos que deben estar incluidos en nuestras depreciaciones). El capital de trabajo necesario para comenzar a trabajar y los activos nominales.

En la Figura 10.19 se determina cuánto será el dinero que se debe gestionar como financiamiento y el que debe invertirse por los socios del proyecto o el inversionista. Los fondos de financiamiento son los que se utilizan para el cálculo de los gastos financieros en la Figura 10.13

\begin{table}[H]
\centering
\scriptsize
\renewcommand{\arraystretch}{1.15}
\setlength{\tabcolsep}{7pt}
\begin{tabular}{|l|r|r|r|}
\hline
\multicolumn{4}{|c|}{\textbf{Plan de Inversión Servicio de Detección de Ofensividad}} \\ \hline
\textbf{ACTIVOS} & \textbf{MONTO} & \textbf{FONDOS PROPIOS} & \textbf{FINANCIAMIENTO} \\ \hline
\textbf{I ACTIVOS FIJOS} & \textbf{22,500.00} & \textbf{4,500.00} & \textbf{18,000.00} \\ \hline
Mobiliario y equipo & 22,500.00 & 4,500.00 & 18,000.00 \\ \hline
 &  &  &  \\[-0.25cm] \hline
\textbf{II CAPITAL DE TRABAJO} & \textbf{-6,945.71} & \textbf{-6,945.71} & \textbf{0.00} \\ \hline
Costos de Producción & 259.81 & 259.81 & 0.00 \\ \hline
Sueldos y carga social A y V & 228.89 & 228.89 & 0.00 \\ \hline
Gastos de venta & 116.64 & 116.64 & 0.00 \\ \hline
Gastos de Administración & 59.64 & 59.64 & 0.00 \\ \hline
Menos \% Venta Contado & -7,610.69 & -7,610.69 & 0.00 \\ \hline
 &  &  &  \\[-0.25cm] \hline
\textbf{III ACTIVOS NOMINALES} & \textbf{3,600.00} & \textbf{3,600.00} & \textbf{0.00} \\ \hline
Constitución & 600.00 & 600.00 & 0.00 \\ \hline
Honorarios & 3,000.00 & 3,000.00 & 0.00 \\ \hline
 &  &  &  \\[-0.25cm] \hline
\textbf{TOTAL} & \textbf{19,154.29} & \textbf{6.03\%} & \textbf{93.97\%} \\ \hline
 &  & \textbf{1,154.29} & \textbf{18,000.00} \\ \hline
\end{tabular}

\vspace{0.2cm}
\textbf{Figura 10.19.} Plan de Inversión \textbf{[Elaboración Propia]}
\end{table}

El cálculo del capital de trabajo dependerá de muchos factores como es el tiempo en que la empresa tardará en comenzar a recibir efectivo propio de la inversión. Para ello tomaremos tantos factores como sean necesarios al proyecto, además restamos el porcentaje de las ventas de contado por ser el dinero que ya estamos recibiendo.
